%!TEX root = rapport.tex

\chapter{Implementation}

In this Chapter we will briefly discuss the implementation of the concepts
presented in the preceding Chapters. The presentation here is only
superficial; for details the actual source code must be consulted.

\section{The Parser}

We have implemented a parser for the grammar specified in Definition
\ref{def:grammar} with a few subtle differences. First of all, the symbol used
for $\lambda$ is a \texttt{\#}. Secondly, all lambda expressions fed into the
parser are run through a preprocessing step that adds all implied parentheses
to the string. This step is included to make the parser definition itself
simpler. The preprocessing step ensures that Conventions \ref{con:assoc} and
\ref{con:extent} are followed.

We used a parser generator to make the parser, since we found it simpler and
more straightforward than to implement it from scratch. There exist numerous
parser generators for Python \cite{pythonParserGeneratorList}. We chose YAPPS2
(Yet Another Python Parser System 2) \cite{YAPPS2} because it is designed such
that it generates human-readable code. YAPPS2 takes a specification of an
LL(1) grammar and writes the Python code for the corresponding parser.

\section{Representation of Lambda Terms}

A lambda term is represented as a directed acyclic graph (DAG). The nodes of the
DAG are encoded in Python as objects. There are three types of nodes: 
abstraction, application and variable:
\begin{description}
	\item[Abstraction nodes] define scopes for bound variables and can be thought
	of as the encoding of the $\lambda$s in the DAG-structure. This node type
	always has one child.
	
	\item[Application nodes] have two children. The first of these
	is either the DAG-representation of some \emph{active} subterm of 
	$M$, or another application node. The second child is the DAG-representation
	of the ``argument'', that is, the lambda term that are to be substituted
	in place for a bound variable. An application node represents a possible
	one-step $\beta$-reduction.
	
	\item[Variable nodes] represent free or bound variables. A bound 
	variable is associated with a scope, represented by an ancestral
	abstraction node. Variable nodes cannot have descendants.
\end{description}

This representation is straightforwardly encoded in Python with a general node
class, \pycode{LambdaNode}, which has a subclass for each node type.
\begin{itemize}
	\item Each \pycode{LambdaNode}-object maintains pointers to its immediate 
	descendants as well as a pointer to its parent node.
	
	\item The abstraction nodes furthermore maintain a list of pointers to the
	variable nodes that represent the variables bound by this particular abstractor.
	These variable nodes must be descendants of the abstraction node, although 
	not necessarily direct descendants.
	
	\item Every abstraction node knows the name of the variable that it binds.
	This is done in order to be able to easily build the list of bound 
	variables after parsing a string representation of a lambda term. When
	the list of variables has been built, this name is not important anymore, 
	but it still serves as a means of producing human-readable output of the DAG.
	
	\item When a term is reduced, the problem is attacked ``top-down'', i.e. we first
	see the abstraction node, and then we have to deal with its variables.
	Therefore there is no need for a doubly linked connection between abstraction
	nodes and variable nodes, and we can safely encode the relationship as a simple
	list of bound nodes in each abstraction node.
\end{itemize}

A term that is not in $\beta$-normal-form has at least one $\beta$-contractum,
and therefore the DAG-representation of that term contains at least one
application node and one abstraction node. More specifically, for each
possible one-step $\beta$-reduction in a term $M$, its DAG will contain one
application node. The first descendant of this node will be an abstraction
node, while the second, ``rightmost'', descendant will be the argument that
will be applied to the function represented by the abstraction node.

%Given this DAG, it is possible to implement the aforementioned operations on lambda
%terms by treewalks. 


% Beskrivelse af repræsentation af reduktionsgrafer
\section{Representation of $\beta$-reduction Graphs}

A reduction graph is represented in Python as an instance of the
\pycode{Graph}-class. A \pycode{Graph}-object contains its nodes, represented
as instances of the \pycode{Node}-class, in three different data structures:
\begin{enumerate}
	\item A dictionary with references to the nodes, where the keys are normalized string 
	representations of the lambda terms represented by the nodes.
	
	\item Another dictionary with references to the same nodes, but keyed on the
	names of the nodes instead of the lambda terms.
	
	\item A list of the nodes. This is included to ensure that the nodes can be
	accessed in a consistent ordering, since the Python function \pycode{dict.items()}
	returns the dictionarys elements in an arbitrary order.
\end{enumerate}
A \pycode{Graph}-object contains exactly one \pycode{Node} representing a
lambda term; if one tries to add an existing node, nothing happens.

A \pycode{Node}-object has a name and a reference to the lambda term that it
represents. It furthermore has a list of its children and a list of its
parents. These lists contain instances of the \pycode{Edge}-class, which
basically just implements a coupling between two \pycode{Node}-objects.

% Beskrivelse af funktionerne der genererer reduktionsgraferne
\section{Generation of $\beta$-reduction Graphs}

In the file \texttt{operations.py}, the functions necessary for the generation
of reduction graphs are implemented. The main function here is
\pycode{reductiongraphiter()} which returns an iterator over the different
``generations'' of the reduction graph. Each generation is a reduction graph
containing everything from the generation immediately before, along with the
contracta of all the redexes of one term that was not reduced in the graph of
the preceding generation. Thus, the longer one iterates, the bigger and more
complete a reduction graph is returned. 

The function takes three arguments: the lambda term in whose reduction graph
we are interested, a start parameter controlling when the iterator should
start emitting graphs (making it possible to return only the ten last graphs)
and an end parameter that controls for how long the computation shall
continue, i.e. how many terms should be reduced.

In the case of reduction graphs of infinite size, the computation only halts
when the given maximum is reached. For finite graphs, the computation halts
either when the maximum is reached or when the graph is fully computed,
depending on the chosen maximum value.

Each time a reduction graph is returned by Pythons \pycode{yield}-statement, a
deep copy operation is performed. This is done in order to ensure that later
stages in the program indeed sees the different generations of graphs as
distinct objects, instead of a collection of references to the same set of
\pycode{Node}s. The built-in copy operation \pycode{copy.deepcopy()} proved to
be too slow, so a custom deep copy operation was implemented.

% Beskrivelse af mulighederne for udvidelse med nye tegnealgoritmer
\section{Drawing Algorithms}

A drawing algorithm is implemented in a class that inherits from
\pycode{Graph}, i.e. it is itself a reduction graph that implements a
``drawing interface'', specifying the methods that start the actual algorithm
that assigns x- and y-coordinates to nodes. To extend the software with new
drawing algorithms, a new class extending the \pycode{Graph}-class and the
\pycode{DrawingAlgorithm}-class must be added. The
\pycode{DrawingAlgorithm}-class merely serves as an interface, guaranteeing
that all required methods are present.


